% !TeX root = ../main.tex
\chapter{Hardware Deployment and Preliminary Experimental Results}
\label{ch:hardware}

The dummy loop of \cref{ch:simulation} accepts the Experiment~A files on
a kinematic plant. This chapter takes the same files through the MRT
bridge of \cref{sec:hw-interfaces}: first with mock
\texttt{ros2\_control} interfaces that echo commands into state, then on
the physical \hero{} platform. The mock still has policy lag, velocity
clamps and the Zenoh boundary; what it does not have is actuator
dynamics or wheel-odometry drift. The gap between the two sessions is
therefore the contribution of the plant and of the state estimate, on
top of a residual that the dummy never sees.

\Cref{sec:hw-deployment} records how the deployment was started
and which sessions are reported. \Cref{sec:hw-challenges} treats the
timing and model problems that show up once the bridge is in the loop.
\Cref{sec:hw-results} compares one mock lap with one hardware lap on the
eight-corner path of \cref{fig:sim-box}, using the metrics of
\cref{tab:sim-metrics}. Both traces are internal metrology
(\cref{sec:sim-metrology}). A single slow lap is not an ISO~9283
test~\cite{ISO9283}, and it is not yet a comparison with a decoupled
baseline.

% ---------------------------------------------------------------------
\section{Deployment on the \hero{} Platform}
\label{sec:hw-deployment}

The topology is that of \cref{sec:distributed-arch}: the
external PC solves the optimal-control problem and publishes the
Experiment~A reference; the on-board IPC evaluates the policy and
commands the drives. Both machines run the project containers on
ROS~2~Jazzy with the Zenoh middleware. The external PC was attached to
the platform subnet by Ethernet for every session below. The wireless
access point of the platform remains a possible join path, and was the
reason Zenoh was chosen in the first place, but it was not used as the
control link here: a cable avoids the jitter that made DDS unusable on
that hop, at the price of tethering a \qty{400}{\kilo\gram} base to a
desk. The box of \cref{fig:sim-box} is small enough that the cable did
not go taut in the recorded laps; it is still a run condition, not a
property of the controller.

Start-up follows the order already required by the reset handshake.
Each machine starts its Zenoh router, the external PC router pointing at
the IPC. The Experiment~A host launch then brings up the \mpc{} node, the
waypoint target node and RViz, with the same task file and waypoint file
the dummy used. Only after the reset service is advertised is the IPC
launch started, so that the bridge can handshake immediately rather than
wait. On the first start after a model change the external PC still
spends about a minute compiling CppAD libraries; later starts reuse the
cache. RViz is checked before anyone calls \texttt{/start}: the landed
box, the orange lead legs, and in particular the lead-out at
$z=\qty{1.15}{\meter}$, have to sit in free space. The approaching gate
of \cref{sec:sim-node} then holds the reference at the lead-in until the
measured \ee{} is within \qty{0.02}{\meter} of that point. Orientation
is not part of the gate, so a residual attitude at entry is admitted
into the first measured edge.

A mock session is the same pair of launches with
\texttt{mock\_arm:=true} and \texttt{mock\_platform:=true} on the IPC
side. The controller manager, the MRT bridge and the three \ocs{}
streams across Zenoh are those of the hardware stack; only the
\texttt{ros2\_control} hardware interfaces are replaced by echo
components~\cite{ros2control}. That is a stricter check than the dummy,
and a weaker one than the robot. It is the last run that can still be
repeated at a desk.

Closed-loop motion on the physical platform was obtained before the
sessions reported here. Those earlier laps are omitted, because the
approaching gate, the segment identifiers and the metric topic were not
yet the pipeline of \cref{ch:simulation}. The traces below are one mock
lap and one hardware lap recorded on 14~August~2026, at the operating
point of \cref{sec:sim-node}: path speed
$\qty{0.05}{\meter\per\second}$, corner dwell $\qty{2.0}{\second}$,
horizon $T=\qty{1.3}{\second}$, solution time window
$T_{\mathrm{sol}}=\qty{1.2}{\second}$, feed-forward SLQ with one
iteration, $\mu_{\mathrm{p}}=16$, $\mu_{\mathrm{o}}=3$, and
$R_{\mathrm{base}}$ pre-scale $5.5$. The solver on the external PC used
eight threads. Host and IPC logs were kept alongside the bags
(\texttt{bag\_host\_0814\_1558} mock,
\texttt{bag\_host\_0814\_1613} hardware). Metrics are computed from the
exported \texttt{ee\_metrics} and solver-diagnostics streams, restricted
to $\texttt{segment\_id}\ge 1$ as in \cref{sec:sim-pose-error}.

% ---------------------------------------------------------------------
\section{Deployment Challenges and Mitigations}
\label{sec:hw-challenges}

The real-time invariant of \cref{eq:rt-invariant} was written for the
everyday file \texttt{task\_hardware.info}, where
$T_{\mathrm{sol}}=\qty{1.0}{\second}$. The Experiment~A file widens that
window to $\qty{1.2}{\second}$ so that a slow solve can still hand the
bridge a policy before the
previous one expires. On the hardware lap the last-solve time had mean
$\qty{227}{\milli\second}$, median $\qty{218}{\milli\second}$ and 95th
percentile $\qty{297}{\milli\second}$; the slowest sample in that stream
was $\qty{353}{\milli\second}$, and the solver's own running maximum
reached $\qty{720}{\milli\second}$. Policies arrived at
$\qty{4.4}{\hz}$. Twice the mean solve, $\qty{0.45}{\second}$, is still
less than $T_{\mathrm{sol}}$, and the bridge did not enter a
policy-exhaustion hold on this lap. The loop therefore ran in real time
in the sense of \cref{eq:rt-invariant}, although policies arrived at
less than half the $\qty{10}{\hz}$ named by \texttt{mpcDesiredFrequency};
on the robot that entry throttles nothing (\cref{sec:hw-interfaces}).

The mock lap was faster on average (mean $\qty{174}{\milli\second}$,
effective $\qty{5.8}{\hz}$) but not safer in the tail. One solve took
$\qty{1.12}{\second}$, which is inside the $\qty{1.2}{\second}$ window
and would have been past the $\qty{1.0}{\second}$ window of the everyday
file. The bridge logged three ``policy nearly exhausted'' warnings on
that lap; the hardware lap logged none. \ocs{} still printed its
``solution time window might be shorter than the MPC delay'' warning
once on the robot and nine times on the mock. The warning compares the
window with an internal delay estimate and is not the same test as a
hold. Widening $T_{\mathrm{sol}}$ is what kept the spike from freezing
the mock plant. It does not make a $\qty{10}{\hz}$ solver out of a
$\qty{4}{\hz}$ one. \Cref{fig:hw-solve-ecdf} puts the two distributions
side by side.

\begin{figure}[htbp]
  \centering
  \begin{tikzpicture}
    \begin{axis}[
        width=0.92\linewidth, height=6.0cm,
        xmode=log,
        log ticks with fixed point,
        xlabel={last-solve time / \si{\milli\second}},
        ylabel={fraction of solves at or below},
        xmin=100, xmax=1300,
        ymin=0, ymax=1.03,
        xtick={100,150,200,300,500,800,1200},
        ytick={0,0.25,0.5,0.75,0.95},
        yticklabels={$0$,$0.25$,$0.50$,$0.75$,$0.95$},
        x tick label style={/pgf/number format/1000 sep={}},
        legend pos=north west,
      ]
      \addplot[RWTHBlau100, thick, const plot, no marks]
        table[x=ms, y=frac] {03_figures/hw_mock_solve_ecdf.dat};
      \addplot[RWTHOrange100, thick, const plot, no marks]
        table[x=ms, y=frac] {03_figures/hw_real_solve_ecdf.dat};
      \addplot[gray, dotted, forget plot, domain=100:1300, samples=2] {0.95};
      \addplot[black, dashed, forget plot] coordinates {(1000,0) (1000,1.03)};
      \addplot[black, densely dashdotted, forget plot]
        coordinates {(1200,0) (1200,1.03)};
      \legend{mock, robot}
      \node[anchor=east, font=\scriptsize] at (axis cs:980,0.58)
        {\qty{1.0}{\second}};
      \node[anchor=east, font=\scriptsize] at (axis cs:1180,0.40)
        {\qty{1.2}{\second}};
    \end{axis}
  \end{tikzpicture}
  \caption[Distribution of last-solve time on the mock and hardware
    sessions]{%
    Empirical distribution of the last-solve time on the two laps, on a
    logarithmic axis; \cref{fig:hw-solve} shows the same two streams
    against wall time. The mock curve sits to the left of the robot curve
    over most of its range---the mock solves faster at the median,
    \qty{163}{\milli\second} against \qty{218}{\milli\second}---but it is
    the mock that carries the tail: its slowest solve is
    \qty{1.12}{\second}, against \qty{353}{\milli\second} on the robot,
    and about one mock solve in a hundred exceeds \qty{400}{\milli\second}
    where the robot stream ends. The dashed line is the solution time
    window of \texttt{task\_hardware.info} and the dash-dotted line the
    wider window of the Experiment~A file; the mock spike falls between
    them, which is why it expired a policy under the everyday setting but
    not under the one actually used (\cref{eq:rt-invariant}).}
  \label{fig:hw-solve-ecdf}
\end{figure}

Because a new policy is already old when it arrives, the input that the
bridge applies is never the one at $\tau=0$ of the horizon.
\Cref{fig:hw-lag} sketches the timing. On the hardware lap the lag at
policy arrival---solve time plus the Zenoh hop---had mean
$\qty{0.24}{\second}$ and 95th percentile $\qty{0.32}{\second}$ after
the first few samples, against a maximum of $\qty{0.36}{\second}$.
Between arrivals the bridge keeps evaluating the same policy, so the lag
at the interpolation site grows until the next message; that quantity
had mean $\qty{0.36}{\second}$. The mock figures were
$\qty{0.18}{\second}$ at arrival and $\qty{0.29}{\second}$ at
interpolation. In both cases the first one to two tenths of a second of
every plan are unused. Diagnostics
\texttt{[policy\_arrival]} and \texttt{[policy\_interp]} were left
enabled at $\qty{0.1}{\second}$ on the IPC launch so that these numbers
could be read from the log rather than inferred from solve time alone.

\begin{figure}[htbp]
  \centering
  \begin{tikzpicture}[
      font=\footnotesize,
      >=Latex,
      lbl/.style={font=\scriptsize, align=center},
    ]
    \draw[->] (0,0) -- (12.2,0) node[right] {$\tau$};
    \foreach \x/\t in {0/0, 2.2/0.24, 3.3/0.36, 11.0/1.3} {
      \draw (\x, 0.12) -- (\x, -0.12);
      \node[below, lbl] at (\x, -0.12) {$\t\,\mathrm{s}$};
    }
    \draw[very thick, RWTHBlau100] (0, 0.85) -- (11.0, 0.85);
    \node[lbl, anchor=west] at (0, 1.25)
      {policy valid on $[t_0,\,t_0+T]$, $T=\qty{1.3}{\second}$};
    \fill[RWTHOrange100] (0, 0.85) circle (2.2pt);
    \node[lbl, above] at (0, 0.95) {$t_0$ snapshot};
    \draw[RWTHOrange100, thick] (0, 0.55) -- (2.2, 0.55);
    \node[lbl, above] at (1.1, 0.55) {solve $+\mathrm{net}$};
    \fill[RWTHBlau100] (2.2, 0.85) circle (2.2pt);
    \node[lbl, above] at (2.2, 0.95) {arrival};
    \draw[thick, dashed, RWTHOrange100] (2.2, 0.85) -- (3.3, 0.85);
    \fill[black] (3.3, 0.85) circle (2.2pt);
    \node[lbl, above] at (3.6, 1.15) {typical evaluation};
    \draw[<->] (0, -0.85) -- (2.2, -0.85);
    \node[lbl, below] at (1.1, -0.85)
      {unused prefix $\approx\qty{0.24}{\second}$};
  \end{tikzpicture}
  \caption[Policy lag on the hardware lap]{%
    Policy lag on the hardware lap. A plan stamped at the snapshot
    $t_0$ becomes available only after the solve and the network hop,
    and is then read at a time that continues to advance until the next
    plan arrives. Times are the measured means of that session, not
    configuration values.}
  \label{fig:hw-lag}
\end{figure}

If that prefix contains the only aggressive correction the solver is
willing to apply, the robot never receives it. The \ocs{} model is a
kinematic integrator: a large velocity at $\tau=0$ really would cancel a
position error in a few tens of milliseconds, and a high
$\mu_{\mathrm{p}}$ encourages exactly that pulse. The recorded policy
topic stores the first input of each plan, not the input at the
interpolation time, so this chapter cannot show the pulse being skipped
on these laps. What it can show is a consistency check. Along-track lag
$t_{\mathrm{lag}}$ of \cref{eq:lag-eff} should be comparable to the
interpolation lag if delay is the whole story. On the mock it is
(mean $t_{\mathrm{lag}}=\qty{0.39}{\second}$ against
$\qty{0.29}{\second}$ of interpolation lag). On the robot it is not
(mean $t_{\mathrm{lag}}=\qty{1.81}{\second}$). Delay explains the mock
residual and does not explain the hardware residual.

Two other failure modes of the bridge were already designed around and
barely appeared in these sessions. A mid-run reset handshake still
replaces the reference with a two-knot hold at the current \ee{}; the
waypoint node re-publishes the same knots at $\qty{1}{\hz}$ on the
\ocs{} time axis so that a lost reference comes back without a step in
the command. The arm rate clamp of \cref{tab:bridge-guards} issued two
warnings on each lap, which is not a dominant trace. Protective stops
and teach-pendant mode errors, which dominated the first bring-up of the
IPC launch, are operational rather than closed-loop problems: the UR
driver has to start in headless remote mode, and the controller stopper
has to keep the joint-trajectory controller in its ignore list, as
already required in \cref{sec:hw-interfaces}. They are mentioned here
only because a session that skips those checks never produces a bag.

The remaining limitation is the model. Mock hardware realises the
integrator: commanded wheel and joint rates appear in the next
observation. The physical base is a \qty{400}{\kilo\gram} platform whose
drive controller enforces its own acceleration
limit~\cite{HenkelRoth2019freeMoVe}, and the UR10e in streaming position
control has a following error that the kinematic model does not
know~\cite{UR10e}. Every re-solve on the robot therefore starts from a
state the previous plan did not expect. That is a different problem from
a slow solver, and \cref{sec:hw-results} is mostly about it.

% ---------------------------------------------------------------------
\section{Preliminary End-Effector Tracking Results}
\label{sec:hw-results}

Both sessions completed the eight measured edges and finished on the
lead-out hold at $z=\qty{1.15}{\meter}$. The mock approaching gate ran
for $\qty{10}{\second}$ from an initial \ee{} error of
$\qty{1.04}{\meter}$; the robot needed $\qty{19}{\second}$ from
$\qty{0.34}{\meter}$, longer than the mock from a third of the
distance, and entered \texttt{RUNNING} at
$\qty{19}{\milli\meter}$ with an orientation error of
$\qty{0.061}{\radian}$ that the gate does not wait for. From the first
sample with $\texttt{segment\_id}\ge 1$ to the last, both laps lasted
$\qty{104}{\second}$, which is the commanded
$\qty{4.50}{\meter}/\qty{0.05}{\meter\per\second}$ plus seven corner
dwells of $\qty{2}{\second}$ (the last corner has no dwell before the
lead-out). Task-completion time $T$ of \cref{tab:sim-metrics} therefore
does not distinguish the two plants. Path shape does.

\Cref{fig:hw-path} shows the commanded and attained \ee{} paths in a
frame attached to the first corner, with $x$ along P1--P2 and $y$ along
the leftward horizontal edge. The vertical edges of the box collapse in
this view; they remain in the error signals. On the mock the attained
path follows the rectangle with a lag that is visible mainly as a
short-cut at the corners. On the robot the attained path is a rounded,
inward-offset copy of the same rectangle, with a consistent bias of
order \qty{10}{\centi\meter} rather than a delay that stays on the
line.

\begin{figure}[htbp]
  \centering
  \begin{subfigure}[t]{0.48\textwidth}
    \centering
    \begin{tikzpicture}
      \begin{axis}[
          width=0.95\linewidth, height=6.1cm,
          xlabel={$x$ / \si{\meter}},
          ylabel={$y$ / \si{\meter}},
          xmin=-0.12, xmax=0.78, ymin=-0.12, ymax=0.95,
          axis equal image,
          legend pos=north west,
        ]
        \addplot[black, thick] table {03_figures/hw_mock_des.dat};
        \addplot[RWTHBlau100, thick] table {03_figures/hw_mock_act.dat};
        \legend{commanded, attained}
      \end{axis}
    \end{tikzpicture}
    \caption{Mock hardware}
    \label{fig:hw-path-mock}
  \end{subfigure}
  \hfill
  \begin{subfigure}[t]{0.48\textwidth}
    \centering
    \begin{tikzpicture}
      \begin{axis}[
          width=0.95\linewidth, height=6.1cm,
          xlabel={$x$ / \si{\meter}},
          ylabel={$y$ / \si{\meter}},
          xmin=-0.12, xmax=0.78, ymin=-0.12, ymax=0.95,
          axis equal image,
          legend pos=north west,
        ]
        \addplot[black, thick] table {03_figures/hw_real_des.dat};
        \addplot[RWTHOrange100, thick] table {03_figures/hw_real_act.dat};
        \legend{commanded, attained}
      \end{axis}
    \end{tikzpicture}
    \caption{Physical robot}
    \label{fig:hw-path-real}
  \end{subfigure}
  \caption[Commanded and attained \ee{} paths, mock versus hardware]{%
    Commanded and attained \ee{} paths on the measured edges, in a
    frame with origin at P1 and $x$ along P1--P2. The two vertical
    edges of \cref{fig:sim-box} project onto the short sides.}
  \label{fig:hw-path}
\end{figure}

\Cref{tab:hw-metrics} collects the scalar scores. Time-aligned position
error $e_{\mathrm{p}}$ on the mock had RMS \qty{27}{\milli\meter} and
maximum \qty{71}{\milli\meter}; on the robot the same quantities were
\qty{130}{\milli\meter} and \qty{204}{\milli\meter}. Mean Hausdorff
$\overline{H}$ rose from \qty{12}{\milli\meter} to
\qty{62}{\milli\meter}, and the 95th percentile of the nearest-neighbour
distances from \qty{23}{\milli\meter} to \qty{110}{\milli\meter}. The
classical $H$ is dominated by a single mock sample at
\qty{71}{\milli\meter} and is the weaker summary, which is why
\cref{sec:sim-path-metrics} asked for $\overline{H}$ and $P_{95}$ as
well. Orientation error grew with the plant too: RMS
\qty{0.056}{\radian} on the mock, \qty{0.27}{\radian} on the robot, with
a hardware maximum of \qty{0.47}{\radian} on later edges. With
$n=1$ cycle, ISO path accuracy $\mathrm{AT}_{p}$ is not reported as a
standard value; $h(B,A)$ on the densely resampled command is the nearest
stand-in and sits at \qty{71}{\milli\meter} (mock) and
\qty{123}{\milli\meter} (robot).

\begin{table}[htbp]
  \centering
  \caption[Mock versus hardware scores on the measured Experiment~A lap]{%
    Mock versus hardware scores on the measured Experiment~A lap
    ($\texttt{segment\_id}\ge 1$). Hausdorff quantities use a
    \qty{1}{\milli\meter} resampling of the commanded polyline, as in
    \cref{sec:sim-path-metrics}. Solve time is the last-solve sample of
    the diagnostics stream; policy lag at arrival is the MRT
    \texttt{policyLagAtArrival} log, excluding the first five samples.}
  \label{tab:hw-metrics}
  \begin{tabularx}{\linewidth}{@{}lXX@{}}
    \toprule
      & Mock hardware & Physical robot \\
    \midrule
    $e_{\mathrm{p}}$ RMS (max)
      & \qty{27}{\milli\meter} (\qty{71}{\milli\meter})
      & \qty{130}{\milli\meter} (\qty{204}{\milli\meter}) \\
    $e_{\mathrm{o}}$ RMS (max)
      & \qty{0.056}{\radian} (\qty{0.13}{\radian})
      & \qty{0.27}{\radian} (\qty{0.47}{\radian}) \\
    $H$ / $\overline{H}$ / $P_{95}$
      & \qty{71}{\milli\meter} / \qty{12}{\milli\meter} /
        \qty{23}{\milli\meter}
      & \qty{135}{\milli\meter} / \qty{62}{\milli\meter} /
        \qty{110}{\milli\meter} \\
    mean $t_{\mathrm{lag}}$
      & \qty{0.39}{\second}
      & \qty{1.81}{\second} \\
    mean solve time (effective rate)
      & \qty{174}{\milli\second} (\qty{5.8}{\hz})
      & \qty{227}{\milli\second} (\qty{4.4}{\hz}) \\
    mean policy lag at arrival
      & \qty{0.18}{\second}
      & \qty{0.24}{\second} \\
    measured lap $T$
      & \qty{104}{\second}
      & \qty{104}{\second} \\
    \bottomrule
  \end{tabularx}
\end{table}

\Cref{fig:hw-error} plots $e_{\mathrm{p}}(t)$ on the measured edges,
with time zero at the first admitted sample. The mock trace stays below
\qty{50}{\milli\meter} except for a short spike at the first corner. The
hardware trace climbs through the first long edge, plateaus around
\qty{15}{\centi\meter} to \qty{20}{\centi\meter} on edges 2--4 and~6, and
only comes back toward \qty{5}{\centi\meter} on the last horizontal
edge. It is not a delayed copy of the mock.

\begin{figure}[htbp]
  \centering
  \begin{tikzpicture}
    \begin{axis}[
        width=0.92\linewidth, height=5.6cm,
        xlabel={time on measured path / \si{\second}},
        ylabel={$e_{\mathrm{p}}$ / \si{\milli\meter}},
        xmin=0, xmax=105, ymin=0, ymax=220,
        legend pos=north west,
      ]
      \addplot[RWTHBlau100, thick] table {03_figures/hw_mock_ep.dat};
      \addplot[RWTHOrange100, thick] table {03_figures/hw_real_ep.dat};
      \legend{mock, robot}
    \end{axis}
  \end{tikzpicture}
  \caption[Time-aligned \ee{} position error on the measured path]{%
    Time-aligned \ee{} position error on the measured path. Time zero is
    the first sample with $\texttt{segment\_id}\ge 1$.}
  \label{fig:hw-error}
\end{figure}

Per-edge numbers make the same split (\cref{tab:hw-segments}). The mock
RMS stays between \qty{10}{\milli\meter} and \qty{41}{\milli\meter} on
every edge. The robot is already at \qty{69}{\milli\meter} RMS on the
opening \qty{0.60}{\meter} edge and peaks at \qty{175}{\milli\meter} on
edge~3. The two vertical edges (4 and~8), which the base cannot help
with, are not the worst hardware edges; the long horizontal edges are.
Along-track error on the robot is positive on seven of the eight edges
and has mean \qty{91}{\milli\meter}, which at
$v_{\mathrm{ref}}=\qty{0.05}{\meter\per\second}$ is the
$\qty{1.81}{\second}$ of $t_{\mathrm{lag}}$ already quoted. A pure
policy-lag prediction $0.36\,\mathrm{s}\times v_{\mathrm{ref}}$ is only
\qty{18}{\milli\meter}. The mock along-track mean, \qty{19}{\milli\meter},
matches that prediction. The extra \qty{70}{\milli\meter} on the robot
is not delay in the MRT bridge.

\begin{table}[htbp]
  \centering
  \caption[Per-edge \ee{} position error on the measured lap]{%
    Per-edge \ee{} position error on the measured lap. Edges 1, 3, 5
    and~7 are the \qty{0.60}{\meter} sides; 2 and~6 are the
    \qty{0.80}{\meter} sides; 4 and~8 are the vertical \qty{0.25}{\meter}
    sides.}
  \label{tab:hw-segments}
  \begin{tabularx}{\linewidth}{@{}lXXXX@{}}
    \toprule
    Edge & Mock RMS & Mock max & Robot RMS & Robot max \\
    \midrule
    1 & \qty{16}{\milli\meter} & \qty{71}{\milli\meter}
      & \qty{69}{\milli\meter} & \qty{80}{\milli\meter} \\
    2 & \qty{41}{\milli\meter} & \qty{53}{\milli\meter}
      & \qty{143}{\milli\meter} & \qty{187}{\milli\meter} \\
    3 & \qty{10}{\milli\meter} & \qty{18}{\milli\meter}
      & \qty{175}{\milli\meter} & \qty{204}{\milli\meter} \\
    4 & \qty{31}{\milli\meter} & \qty{45}{\milli\meter}
      & \qty{145}{\milli\meter} & \qty{174}{\milli\meter} \\
    5 & \qty{22}{\milli\meter} & \qty{26}{\milli\meter}
      & \qty{98}{\milli\meter} & \qty{129}{\milli\meter} \\
    6 & \qty{30}{\milli\meter} & \qty{35}{\milli\meter}
      & \qty{169}{\milli\meter} & \qty{202}{\milli\meter} \\
    7 & \qty{18}{\milli\meter} & \qty{21}{\milli\meter}
      & \qty{60}{\milli\meter} & \qty{140}{\milli\meter} \\
    8 & \qty{39}{\milli\meter} & \qty{45}{\milli\meter}
      & \qty{111}{\milli\meter} & \qty{160}{\milli\meter} \\
    \bottomrule
  \end{tabularx}
\end{table}

Corner dwells isolate the same point without mixing it into a moving
reference. On the mock, the last $\qty{2}{\second}$ of edges 1--7---the
dwell windows---fall to between \qty{5}{\milli\meter} and
\qty{25}{\milli\meter}. On the robot the same windows remain between
\qty{53}{\milli\meter} and \qty{168}{\milli\meter} and do not collapse.
Once the command has stopped, a lag of $\qty{0.36}{\second}$ should
have died out. What remains is a steady offset that the receding-horizon
loop does not remove at this operating point.

Solve time over the two sessions is shown in \cref{fig:hw-solve} only to
close the timing discussion. After the first seconds both traces sit
near their means; the mock spike to $\qty{1.12}{\second}$ is the event
that produced the exhaustion warnings. Neither curve tracks
$e_{\mathrm{p}}(t)$. A slow solver would have frozen the robot through
the hold logic of \cref{tab:bridge-guards}. That did not happen. The
tracking gap is left by the plant and by the lag acting on a model that
assumes the plant is an integrator.

\begin{figure}[htbp]
  \centering
  \begin{tikzpicture}
    \begin{axis}[
        width=0.92\linewidth, height=5.2cm,
        xlabel={wall time from first diagnostic / \si{\second}},
        ylabel={last solve time / \si{\second}},
        xmin=0, xmax=210, ymin=0, ymax=1.25,
        legend pos=north east,
      ]
      \addplot[RWTHBlau100, thick] table {03_figures/hw_mock_solve.dat};
      \addplot[RWTHOrange100, thick] table {03_figures/hw_real_solve.dat};
      \addplot[black, dashed, domain=0:210] {0.1};
      \addplot[gray, dashed, domain=0:210] {1.2};
      \legend{mock, robot, $\qty{10}{\hz}$ period, $T_{\mathrm{sol}}$}
    \end{axis}
  \end{tikzpicture}
  \caption[Last-solve time on the mock and hardware sessions]{%
    Last-solve time on the mock and hardware sessions. The lower dashed
    line is the $\qty{0.1}{\second}$ period of \texttt{mpcDesiredFrequency},
    which the bridge uses only as its warning threshold; the upper line is
    $T_{\mathrm{sol}}=\qty{1.2}{\second}$.}
  \label{fig:hw-solve}
\end{figure}

Two caveats bound what these numbers may be used for. First, both
desired and attained poses are reconstructed from the \mpc{} observation,
so wheel-odometry drift is shared by the command frame and the measured
\ee{}. On the robot the box was landed at a base pose near
$(23.2,\,9.0)$ in the odometry frame, which is the accumulated motion of
the platform before this lap, not a localisation in the hall. Closed-loop
drift of the attained path relative to its own start was
\qty{15}{\centi\meter} on the robot and \qty{11}{\centi\meter} on the
mock, most of it in $z$ at the last vertical edge rather than a planar
loop-closure error. None of this is an ISO accuracy statement; an
external \ee{} pose is still missing. Second, each plant was run once, at
one speed, without a decoupled baseline. Repeatability $\mathrm{RT}_{p}$,
the three speed levels of ISO Table~3, and a Nav2 plus MoveIt comparison
are outside the present
data~\cite{ISO9283,Macenski2020Nav2,moveit2}.

What the pair of laps does show is narrower. The Experiment~A files that
the dummy accepted can be executed by the hardware stack from approaching
gate to lead-out hold. Mock hardware, which shares the bridge and the
lag with the robot, tracks at a few centimetres and at an along-track
delay consistent with policy lag. The physical platform, with the same
files, tracks at one to two decimetres and does not catch the reference
during dwell. The next improvements that would change those numbers are
a shorter solve (so that $u(\tau=0)$ is no longer in the past), a model
that knows the base cannot step in velocity, and an external pose of the
\texttt{tool0} frame.
